\chapter[Linear algebra and geometry]{Vector spaces, matrices, and high-dimensional geometry}
Let $X\in\R^{T\times N}$ collect $N$ neural variables across $T$
observations. Its columns may be neurons, channels, or features;
centering across rows changes the statistical quantity being modeled.
Linear algebra does not specify how trials should be aligned or how
observations should be validated.

\entry{25}{Rank--nullity theorem}
{For linear $A:V\to W$ with $\dim V<\infty$,
$\dim\ker A+\rank A=\dim V$.}
{Count identifiable and unidentifiable latent directions.}
\entry{26}{Orthogonal projection theorem}
{In a Hilbert space every element has a unique closest
point in a closed subspace; the residual is orthogonal to that subspace.}
{Derive least squares and orthogonal signal decompositions.}
\entry{27}{Spectral theorem}
{A real symmetric matrix admits an orthonormal eigenbasis and real
eigenvalues; a covariance matrix has nonnegative eigenvalues.}
{Explain the eigenvectors underlying population PCA.}
\entry{28}{Singular value decomposition}
{Every real matrix $X$ admits $X=U\Sigma V^\top$
with orthogonal factors and nonnegative singular values.}
{Implement low-rank population representations.}
\entry{29}{Eckart--Young--Mirsky theorem}
{Truncating SVD to rank $k$ gives an optimal rank-$k$ approximation
in Frobenius and spectral norms, more generally unitarily invariant norms.}
{Give a reconstruction criterion for reduced-dimensional activity.}
\entry{30}{Courant--Fischer min--max theorem}
{Eigenvalues of a symmetric matrix have variational characterizations
over subspaces via Rayleigh quotients.}
{Derive why sequential PCA directions maximize remaining variance.}
\entry{31}{Gauss--Markov theorem}
{With a full-column-rank linear model, zero-mean errors
with common variance and no correlation, OLS is BLUE.}
{Clarify which optimality claim linear decoders can and cannot make.}
\entry{32}{Schur complement identities}
{For an invertible block, Schur complements characterize
positive definiteness and block inverses under the stated conditions.}
{Compute conditional Gaussian covariance and latent uncertainty.}
\entry{33}{Weyl eigenvalue inequalities}
{For Hermitian $A,E$, sorted eigenvalues of $A+E$ differ
from those of $A$ by at most $\|E\|_2$.}
{Control eigenspectrum changes caused by covariance perturbations.}
\entry{34}{Davis--Kahan sin-theta theorem}
{For symmetric matrices with separated spectral clusters,
their invariant subspaces have angle bounds proportional
to perturbation size divided by a suitable eigengap.}
{Assess when estimated neural subspaces are stable.}
\entry{35}{Sherman--Morrison--Woodbury identity}
{For invertible $A,C$ and required inverse,
$(A+UCV)^{-1}=A^{-1}-A^{-1}U(C^{-1}+VA^{-1}U)^{-1}VA^{-1}$.}
{Speed covariance updates when their rank is small.}
\entry{36}{Matrix determinant lemma}
{For invertible $A,C$,
$\det(A+UCV)=\det(C^{-1}+VA^{-1}U)\det(C)\det(A)$.}
{Accelerate Gaussian latent-variable likelihood evaluations.}
\entry{37}{Mercer's theorem}
{For a continuous positive-semidefinite kernel on a compact
domain with a suitable finite reference measure,
an orthogonal eigenfunction expansion exists under its hypotheses.}
{Understand kernel feature representations in GP modeling.}
\entry{38}{Karhunen--Loeve theorem}
{A second-order process admits a mean-square orthogonal expansion
in covariance eigenfunctions under suitable regularity.}
{Link temporal covariance to latent trajectories and function-space PCA.}
\entry{39}{Kruskal's uniqueness criterion for CP tensors}
{A canonical polyadic tensor decomposition is essentially unique
if its factor Kruskal ranks meet an appropriate sufficient bound,
e.g.\ $k_A+k_B+k_C\geq 2R+2$ for three factors.}
{Reason about identifiability in trial-by-time-by-neuron decompositions.}

\begin{definition}
For covariance $C\succeq0$ with eigenvalues $\lambda_i$,
the participation ratio is
\[
D_{\mathrm{PR}}=\frac{(\sum_i\lambda_i)^2}{\sum_i\lambda_i^2}.
\]
It measures effective spectral dimensionality, not the intrinsic
manifold dimension without further assumptions.
\end{definition}
\begin{caution}
Truncated SVD solves a matrix approximation problem; recovering
biological latent variables is an additional inference problem.
Coherent noise, trial structure, and preprocessing can dominate
the covariance eigenspectrum.
\end{caution}
% --- expanded edition (October 2026) ---
\section{Proofs of selected matrix results}
\begin{proof}[Proof of Result 31 (Gauss--Markov)]
Let $y=X\beta+\varepsilon$ with $X\in\R^{T\times p}$ of full column rank,
$\E\varepsilon=0$ and $\Cov\varepsilon=\sigma^2I$. The OLS estimator is
$\hat\beta=(X^\top X)^{-1}X^\top y$, with $\Cov\hat\beta=\sigma^2(X^\top X)^{-1}$.
Any linear estimator is $\tilde\beta=Ay$; unbiasedness for every $\beta$
requires $AX=I$. Write $A=(X^\top X)^{-1}X^\top+D$; then $DX=0$, and
\[
\Cov\tilde\beta=\sigma^2AA^\top=\sigma^2(X^\top X)^{-1}+\sigma^2DD^\top,
\]
because the cross terms contain $DX=0$. Since $DD^\top\succeq0$,
$\Cov\tilde\beta-\Cov\hat\beta\succeq0$.
\end{proof}
The theorem's hypotheses are restrictive in neural data. Binned spike
counts are heteroskedastic (their variance grows with their mean) and
time bins are autocorrelated; under either violation OLS remains
unbiased but is no longer best, and its textbook standard errors are wrong.
Weighted or generalized least squares, with $\Cov\varepsilon=\sigma^2\Omega$,
restores optimality by whitening with $\Omega^{-1/2}$.

\begin{proof}[Proof of Result 32 (Schur complement)]
Let $M=\begin{psmallmatrix}A&B\\B^\top&C\end{psmallmatrix}$ be symmetric
with $C\succ0$, and $S=A-BC^{-1}B^\top$. The block factorization
\[
M=\begin{pmatrix}I&BC^{-1}\\0&I\end{pmatrix}
\begin{pmatrix}S&0\\0&C\end{pmatrix}
\begin{pmatrix}I&0\\C^{-1}B^\top&I\end{pmatrix}
\]
is verified by multiplication. The outer factors are invertible and
transposes of each other, so $M$ is congruent to $\diag(S,C)$.
Hence $M\succ0$ iff $S\succ0$ (given $C\succ0$), $\det M=\det S\det C$,
and inverting the three factors gives the block inverse, whose top-left
block is $S^{-1}$.
\end{proof}
In Gaussian conditioning (Chapter~\ref{ch:gaussian-toolkit}), $S$ is
exactly the conditional covariance. The determinant identity
$\det M=\det S\det C$ is the chain rule of Gaussian likelihoods:
$\log p(x,y)=\log p(x\mid y)+\log p(y)$.

\begin{proof}[Proof of Result 35 (Woodbury)]
Multiply $A+UCV$ by the proposed inverse. The product is
\[
I-U(C^{-1}+VA^{-1}U)^{-1}VA^{-1}+UCVA^{-1}-UCVA^{-1}U(C^{-1}+VA^{-1}U)^{-1}VA^{-1}.
\]
Collect the terms ending in $(C^{-1}+VA^{-1}U)^{-1}VA^{-1}$:
their coefficient is $-U-UCVA^{-1}U=-UC(C^{-1}+VA^{-1}U)$. Thus those
terms equal $-UCVA^{-1}$, which cancels the remaining $+UCVA^{-1}$.
\end{proof}
For factor analysis with $\Sigma=\Psi+\Lambda\Lambda^\top$, $\Psi$
diagonal $N\times N$ and $\Lambda$ of rank $k\ll N$, Woodbury reduces
$\Sigma^{-1}$ to the inversion of a $k\times k$ matrix:
$\Sigma^{-1}=\Psi^{-1}-\Psi^{-1}\Lambda(I_k+\Lambda^\top\Psi^{-1}\Lambda)^{-1}\Lambda^\top\Psi^{-1}$.
For $N=1000$ neurons and $k=10$ factors the cost drops from $O(N^3)$
to $O(Nk^2)$.

\begin{proof}[Proof of Result 33 (Weyl)]
By Courant--Fischer (Result 30),
\[
\lambda_k(A+E)=\max_{\dim W=k}\ \min_{x\in W,\ \|x\|=1}x^\top(A+E)x.
\]
For every unit vector $x$,
\[
x^\top Ax-\|E\|_2\le x^\top(A+E)x\le x^\top Ax+\|E\|_2.
\]
Taking min over $W$ then max over subspaces preserves the inequalities.
\end{proof}

\section{Subspace stability in practice}
Result 34 (Davis--Kahan) has a usable form: if $\hat\Sigma=\Sigma+E$
and the top-$k$ eigenspace of $\Sigma$ is separated from the rest by
the gap $\delta=\lambda_k-\lambda_{k+1}>0$, then the principal angles
$\Theta$ between the true and estimated top-$k$ subspaces satisfy
\[
\|\sin\Theta\|_2\le\frac{2\|E\|_2}{\delta}.
\]
For a sample covariance of $T$ observations of $N$ roughly Gaussian
variables, $\|E\|_2$ is of order $\|\Sigma\|_2(\sqrt{N/T}+N/T)$. Two
consequences follow. First, principal subspaces are well estimated only
when the eigengap exceeds the sampling fluctuation. Second, individual
components inside a cluster of nearly equal eigenvalues are not
estimable, while the subspace spanned by the whole cluster may be.

\begin{example}[Split-half reliability of PCs]
Split trials into two halves, compute PCA on each, and measure the
principal angles between the top-$k$ subspaces. Small angles for the
subspace with large angles for individual PCs indicate rotation within
a near-degenerate cluster: report the subspace, not the components.
This is an empirical counterpart of the Davis--Kahan bound.
\end{example}

\section{Karhunen--Lo\`eve expansions of two standard processes}
Result 38 states that a mean-square continuous process on $[0,T]$ with
covariance $K(s,t)$ admits $X(t)=\sum_j\sqrt{\lambda_j}\,\xi_j\,e_j(t)$
with uncorrelated unit-variance $\xi_j$, where $(\lambda_j,e_j)$ are
eigenpairs of the covariance operator.
\begin{example}[Brownian motion]
For $K(s,t)=\min(s,t)$ on $[0,1]$, the eigenproblem
$\int_0^1\min(s,t)e(s)\,ds=\lambda e(t)$ differentiates twice into
$-e''=e/\lambda$ with $e(0)=0$, $e'(1)=0$. Hence
$e_j(t)=\sqrt2\sin((j-\tfrac12)\pi t)$ and
$\lambda_j=1/((j-\tfrac12)^2\pi^2)$. The eigenvalues decay as $j^{-2}$:
most of the variance lies in a few slow components. A random walk
analysed by PCA across time therefore produces smooth, sinusoid-like
``components'' that reflect only its covariance, not oscillations.
\end{example}
This example explains a well-known pitfall: PCA of smoothly drifting
population activity yields leading components resembling sines and
cosines of increasing frequency, which can be mistaken for rotational
or oscillatory dynamics.

\section*{Exercises for this chapter}
\begin{exercise}
Show that the OLS residual $y-X\hat\beta$ satisfies $X^\top(y-X\hat\beta)=0$
and interpret this as Result 26.
\end{exercise}
\begin{exercise}
Prove the matrix determinant lemma (Result 36) from the Schur complement
determinant identity applied to $\begin{psmallmatrix}A&-U\\V&C^{-1}\end{psmallmatrix}$.
\end{exercise}
\begin{exercise}
For $\Sigma=\diag(3,2.9,1)$ and a perturbation of norm $0.2$, compare
the Davis--Kahan bound for the top-one and top-two subspaces.
\end{exercise}
\begin{exercise}
Verify that $e_j(t)=\sqrt2\sin((j-\tfrac12)\pi t)$ are orthonormal on $[0,1]$.
\end{exercise}
